\documentclass[../../../../SoftwareRequirementsSpecification.tex]{subfiles}
\begin{document}
% Richiede le macro definite in src/macros/useCaseMacros.tex
% (incluse nel preambolo di SoftwareRequirementsSpecification.tex)

\ucstart
\begin{xltabular}{\textwidth}{|l|X|}
  \hline
  \ucid{PT-17} \\ \hline
  \ucname{Cessazione Rapporto con un Atleta} \\ \hline
  \ucactor{PT} \\ \hline
  \ucdescription{Il PT termina il rapporto con un atleta che non
    segue più.} \\ \hline
  \ucprecond{Il PT deve essere autenticato. Il PT deve seguire almeno
    un atleta.} \\ \hline
  \ucpostcond{Il sistema imposta la data di fine sul rapporto tra il
    PT e l'atleta. Il rapporto non è più attivo. L'atleta non compare
    più nell'elenco degli atleti del PT.} \\ \hline
  \ucmainflow{
    \item Il PT apre la pagina di dettaglio di un atleta (caso d'uso
      PT-08).
    \item Il PT preme il pulsante "Termina Rapporto".
    \item Il sistema chiede conferma. Il sistema avvisa che
      l'operazione non è reversibile e che il PT non potrà più
      accedere ai dati dell'atleta.
    \item Il PT preme il pulsante "Conferma".
    \item Il sistema imposta la data di fine del rapporto alla data
      odierna.
    \item Il sistema mostra un messaggio di conferma e mostra l'elenco
      aggiornato degli atleti del PT.
  } \\ \hline
  \ucaltflow{
    \textbf{4a. Il PT annulla} \newline
    - Il PT preme il pulsante "Annulla" invece di confermare. \newline
    - Il sistema non modifica nulla. \newline
    - Il flusso termina.
  } \\ \hline
  \ucnote{
    \textbf{Nota 1:} il sistema non elimina l'account dell'atleta.
    L'atleta può ancora accedere al sistema. L'atleta vede le schede
    confermate, gli appuntamenti e le palestre (casi d'uso A-03, A-05
    e A-06). L'atleta non può più richiedere appuntamenti, perché non
    ha un PT che lo segue (caso d'uso A-02). \newline
    \textbf{Nota 2:} il sistema non modifica le assegnazioni, gli
    appuntamenti e le Iscrizioni dell'atleta. Le assegnazioni in
    bozza restano incomplete: il PT non può più confermarle. Il PT
    può ancora rispondere alle richieste di appuntamento in attesa
    (caso d'uso PT-06). \newline
    \textbf{Nota 3:} la specifica non prevede un caso d'uso per
    assegnare l'atleta a un altro PT.
  } \\ \hline
\end{xltabular}
\end{document}
